\scanpage{8}\mspage{5}
\begingroup
\makeatletter\tagsleft@true\makeatother
The proof is straightforward. Let $z=\dfrac{x}{-y}$, $w=\dfrac{1}{-y}$, so that $x=\dfrac{z}{w}$ and $y=-\dfrac{1}{w}$. Then in terms of $w$ and $z$ the equation for $E$ is
\[
 w(1-a_1z-a_3w)=z^3+a_2z^2w+a_4zw^2+a_6w^3.
 \tag*{(4.4)}
\]
This shows that we have
\begingroup\scriptsize
\[
\begin{aligned}
w&=z^3+a_1z^4+(a_1^2+a_2)z^5+(a_1^3+2a_1a_2+a_3)z^6+
 (a_1^4+3a_1^2a_2+3a_1a_3+a_2^2+a_4)z^7+\ldots\\
 &=z^3+A_1z^4+A_2z^5+\ldots
\end{aligned}
\tag*{(4.5)}
\]
\endgroup
where $A_\nu$ is a polynomial of weight $\nu$ in the $a_i$ with positive integral coeffs Hence the expansions of $y=-\dfrac{1}{w}$ and $x=-zy$ in terms of $z$ have coefficients in $\mathbb{Z}[a_1,a_2,a_3,a_4,a_6]$. Now
\[
 \frac{\omega}{d\deleted{\unsure{t}}z}=
 \frac{dx/d\deleted{\unsure{t}}z}{2y+a_1x+a_3}=
 \frac{-2\deleted{\unsure{t}}z^{-3}+\ldotp}{-2\deleted{\unsure{t}}z^{-3}+\ldotp}=
 \frac{dy/d\deleted{\unsure{t}}z}{3x^2+2a_2x+a_4-a_1y}=
 \frac{3z^{-4}+\ldotp}{3z^{-4}+\ldots}
\]
has coefficients in $\mathbb{Z}[\frac{1}{2},a_1,\ldots,a_6]$ but also in $\mathbb{Z}[\frac{1}{3},a_1,\ ,a_6]$, hence in $\mathbb{Z}[a_1,\ ,a_6]$. As for the group law, if $z_1$ and $z_2\in\mathfrak{p}$, then the line joining the points $(z_1,w_1)$, $(z_1,w_2)$ in the $(z,w)$-plane has slope $\in\mathfrak{p}^2$, because
\[
 \frac{w_2-w_1}{z_2-z_1}=
 \frac{z_2^3-z_1^3}{z_2-z_1}
 +A_1\frac{z_2^4-z_1^4}{z_2-z_1}+
\]
\margin{\unsure{with|the} $A_i$ as above in (\unsure{4}.5)}
Call this slope $\lambda=\lambda(z_1,z_2)=z_2^2+z_1z_2+z_1^2+A_1(z_2^3+z_2^1z_1+z_1z_1^2+z_1^3)+\unsure{\ldotp\,\ldotp|\ldots}$. Put $\nu=\nu(z_1,z_2)=w_i-\lambda z_i$ $(i=1,2)$. Substituting $w=\lambda(z_1,z_2)z+\nu(z_1,z_2)z$ in (4.4) we find a cubic in $z$ with roots $w_1$ and $w_2$. Looking at the sum of the roots, one sees that the third root $\deleted{\unsure{t}}z_3$ is expressed as a power series in $z_1,z_2$ with coeffs in $R$. In fact,
\[
 z_1+z_2+z_3=
 \frac{a_1\lambda+a_3\lambda^2-a_2\nu-2a_4\lambda\nu-3a_6\lambda^2\nu}
 {1+a_2\lambda+a_4\lambda^2+a_6\lambda^3}
 \tag*{(4.6)}
\]
\margin{(sauf erreur)}
Thus we have the `canonical filtration''
\[
 E(K)\supset E_0(K)\supset E_1(K)\supset E_2(K)\supset\ldots
 \qquad\supset(O)=\bigcap_{m=1}^{\infty}E_m(K)
 \tag*{(4.7)}
\]
with quotients
\begingroup\small
\[
 E(K)\big/E_0(K)\quad\text{finite},\qquad
 E_0\big/E_1\mathrel{\deleted{\unsure{\subset}}\hookrightarrow}\widetilde E_0(k),
 \quad\text{and}\quad E_m\big/E_{m+1}\hookrightarrow k^+
 \quad\text{for }m\geqslant1.
\]
\endgroup
Of course the inclusions $\hookrightarrow$ are bijections $\overset{\sim}{\to}$ if $R$ is complete.
\endgroup

% archive sheet number: 7
